\documentclass[11pt]{article}
\usepackage[T1]{fontenc}
\usepackage{lmodern,amsmath,amssymb,amsthm,mathtools,microtype}
\usepackage[letterpaper,margin=1in]{geometry}
\usepackage[hidelinks]{hyperref}
\setlength{\emergencystretch}{2em}
\newtheorem{theorem}{Theorem}
\newtheorem{lemma}[theorem]{Lemma}
\newtheorem{proposition}[theorem]{Proposition}
\theoremstyle{remark}\newtheorem{remark}[theorem]{Remark}
\newcommand{\R}{\mathbb R}
\newcommand{\E}{\mathbb E}
\newcommand{\ind}{\mathbf 1}
\newcommand{\norm}[1]{\lVert#1\rVert}
\newcommand{\dd}{\,\mathrm d}
\title{Sharp logarithmic corrections under\linebreak stretched exponential target bounds}
\author{}
\date{Supplementary mathematical note\quad 7 October 2026}
\begin{document}
\maketitle
\begin{abstract}
The stretched-exponential upper bound in manuscript 007 has logarithmic
power $2/(3q)$ at Wasserstein exponent $1/3$. Reweighting the explicit
two-atom constructions of manuscript 001 v3 shows that this logarithmic
power is necessary for a fixed uniform cube source and for a fixed
cube-truncated Gaussian source. The latter is strongly log-concave with
a hard convex boundary. In contrast, the full Gaussian source has a
pure one-third bound on the same target class. We give the complete
lower-bound calculation, including the source-cell masses, direct
optimality of the maps, optimal target coupling, and logarithmic
asymptotics.
\end{abstract}

\paragraph{Scope and dependencies.}
This note combines existing proved parts of two research drafts. The
two-atom geometries are those of manuscript 001 v3, Theorem 1.7,
Section 11 and Section 12.2~\cite{M001}. The upper estimate is manuscript
007 v1, Corollary 4.1, using its Theorem 3.1 and Corollary 3.2~\cite{M007}.
The lower bounds below are proved directly and do not use the
radius-independent potential-stability theorem. The source-specific
upper bounds use the potential theorems proved in manuscript 001 v3,
Theorems 1.1 and 1.4; this note does not independently revalidate those
theorems. No novelty, publication-priority,
external-peer-review, or journal-level claim is made. Neither frozen
manuscript is changed by this note.

\section{Statement and common calculation}
For $q,K>0$, define the pair class
\[
 \mathcal E_{q,K}=
 \left\{(\mu,\nu):
 \int e^{(|y|/K)^q}\dd\mu(y)
 +\int e^{(|y|/K)^q}\dd\nu(y)\le4\right\}.
\]
Every member has finite moments of every positive order. For an
absolutely continuous source $\rho$, let $T_\mu$ be its quadratic
Brenier map to $\mu$, and write
\[
 w=W_2(\mu,\nu),\qquad
 X_\rho(\mu,\nu)=\norm{T_\mu-T_\nu}_{L^2(\rho)},
\]
\[
 \Phi_{q,K}(w)=K^{2/3}w^{1/3}
             \bigl(1+\log(K/w)\bigr)^{2/(3q)},
 \qquad 0<w\le K.
\]

\begin{theorem}\label{thm:main}
Fix $d\ge2$ and either of the following sources, independently of
$q,K$ and of the target parameter:
\begin{enumerate}
\item the uniform probability measure on $[-1,1]^d$;
\item the standard Gaussian conditioned on $[-1,1]^d$.
\end{enumerate}
For every $q,K>0$, there are pairs $(\mu_a,\nu_a)\in\mathcal E_{q,K}$,
each target having exactly two atoms, such that $w_a\to0$ as
$a\downarrow0$, and
\[
 \frac{X_\rho(\mu_a,\nu_a)}{\Phi_{q,K}(w_a)}
       \longrightarrow c_{\rho,q}\in(0,\infty).
\]
In particular, no uniform estimate with Wasserstein power $1/3$ and
logarithmic power $\gamma<2/(3q)$ holds on this class. Nor can
$\Phi_{q,K}$ be replaced by a function that is $o(\Phi_{q,K})$ as
$w\downarrow0$.

For these sources, manuscript 007, Corollary 4.1, gives the matching
upper bound $X_\rho\le C_{\rho,d,q}\Phi_{q,K}(w)$ for
$(\mu,\nu)\in\mathcal E_{q,K}$ and $0<w\le K$, using the explicitly
cited potential theorems proved in manuscript 001.
\end{theorem}

The second source has density
\[
 r(x)=Z_d^{-1}e^{-|x|^2/2}\ind_{[-1,1]^d}(x).
\]
It is $1$-strongly log-concave in the extended-potential sense:
$r=Z_d^{-1}e^{-|x|^2/2-W}$ with the proper lower-semicontinuous convex
function $W=I_{[-1,1]^d}$, equal to zero on the cube and $+\infty$
outside. Consequently, the necessity of the logarithm holds even for
one fixed strongly log-concave source with a hard boundary.

\begin{lemma}[Equal-mass two-atom formulas]\label{lem:atoms}
Let $0<m<1/2$, $R>0$, and let $e,n$ be unit vectors with
$e\cdot n>0$. Put $\eta=|e-n|$ and
\[
 \mu=(1-m)\delta_0+m\delta_{Re},\qquad
 \nu=(1-m)\delta_0+m\delta_{Rn}.
\]
Then
\begin{equation}\label{eq:Wexact}
 W_2(\mu,\nu)^2=mR^2\eta^2.
\end{equation}
If their Brenier maps from $\rho$ have the form
$T_\mu=Re\ind_A$ and $T_\nu=Rn\ind_B$, where
$\rho(A)=\rho(B)=m$, and $D=\rho(A\mathbin\triangle B)$, then
\begin{equation}\label{eq:Xexact}
 X_\rho(\mu,\nu)^2
 =R^2\bigl[D+\eta^2(m-D/2)\bigr].
\end{equation}
Finally, taking $R=K[\log(1/m)]^{1/q}$ puts the pair in
$\mathcal E_{q,K}$, with combined exponential moment exactly $4-2m$.
\end{lemma}
\begin{proof}
A coupling is determined by the mass $s\in[0,m]$ sent from $Re$ to
$Rn$. Its other three entries, in the natural two-by-two table, are
$m-s$, $m-s$, and $1-2m+s$. Its cost is
\[
 R^2\bigl[2(m-s)+s\eta^2\bigr]
 =2R^2\bigl[m-s(e\cdot n)\bigr].
\]
Since $e\cdot n>0$, this is minimized at $s=m$, proving
\eqref{eq:Wexact}. The map discrepancy has squared magnitude $R^2$
on $A\mathbin\triangle B$, $R^2\eta^2$ on $A\cap B$, and zero
elsewhere. As $\rho(A\cap B)=m-D/2$, this proves
\eqref{eq:Xexact}. For either target, the exponential moment is
\[
 1-m+m\exp\bigl((R/K)^q\bigr)
 =1-m+m\exp(\log(1/m))=2-m.
\]
\end{proof}

\paragraph{Direct optimality of the ramp maps.}
Both constructions below use finite convex functions that are maxima
of two affine functions, with slopes exactly the two desired target
atoms. To verify optimality without invoking an imported stability
result, write such a function as
$u(x)=\max_{y\in\{0,y_1\}}(x\cdot y+c_y)$. Then
\[
 x\cdot y+c_y\le u(x)
\]
for both target atoms, with equality for the selected gradient slope.
Integrating against any coupling to the same target bounds
$\int x\cdot y$ by the fixed number
$\int u\dd\rho-\sum_y c_y\mu(\{y\})$. The gradient coupling attains
this number. Since both marginal second moments are fixed, it
minimizes quadratic cost. The separating hyperplane is source-null,
so equality selects the unique gradient slope almost everywhere,
also proving uniqueness of the optimal map. All integrals here are
finite, because each source is compact and each individual target
is finite.

\section{Exact uniform-cube construction}
Let $\rho$ be the uniform probability on $[-1,1]^d$. For
$0<a<1/2$, set
\[
 m=\frac a2,\qquad
 n_a=\frac{e_1+a e_2}{\sqrt{1+a^2}},\qquad
 \eta_a=|e_1-n_a|,\qquad
 R_a=K[\log(2/a)]^{1/q}.
\]
Define the rare source cells
\[
 A_a=\{x_1>1-a\},\qquad
 B_a=\{x_1+a x_2>1-a\}.
\]
Conditional on $x_2=z\in[-1,1]$, their horizontal widths inside the
cube are $a$ and $a(1+z)$. Both are in $[0,2]$. The two intervals
share their right endpoint and hence are nested. If $Z$ is uniform
on $[-1,1]$, their source probabilities are therefore exactly
\begin{align}
 \rho(A_a)=\rho(B_a)&=a/2,\label{eq:cube-mass}\\
 D_a=\rho(A_a\mathbin\triangle B_a)
   &=\frac a2\E|Z|=\frac a4,\label{eq:cube-diff}\\
 \rho(A_a\cap B_a)
   &=\frac a2\E\min(1,1+Z)=\frac{3a}{8}.
   \label{eq:cube-intersection}
\end{align}
Here the factor $1/2$ in each horizontal probability is the normalized
one-dimensional cube density. The remaining $d-2$ coordinates
integrate out.

The convex potentials
\[
 u_a(x)=R_a(x_1-1+a)_+,\qquad
 v_a(x)=\frac{R_a}{\sqrt{1+a^2}}(x_1+a x_2-1+a)_+
\]
have gradients $R_a e_1\ind_{A_a}$ and $R_a n_a\ind_{B_a}$.
By \eqref{eq:cube-mass}, their laws are the two-atom targets in
Lemma~\ref{lem:atoms}. The direct argument above proves that these
are their Brenier maps. The lemma now gives the exact identities
\begin{equation}\label{eq:cube-exact}
 w_a^2=mR_a^2\eta_a^2,\qquad
 X_a^2=mR_a^2\left(\frac12+\frac34\eta_a^2\right),\qquad
 \eta_a^2=2-\frac2{\sqrt{1+a^2}}\sim a^2.
\end{equation}
In particular, the combined exponential moment is $4-a\le4$.

Put $\ell_a=\log(2/a)$. Since every fixed power of $\ell_a$ grows
more slowly than every inverse positive power of $a$,
\[
 w_a\sim\frac K{\sqrt2}a^{3/2}\ell_a^{1/q}\longrightarrow0,
 \qquad
 X_a\sim\frac K2 a^{1/2}\ell_a^{1/q}.
\]
Moreover
\[
 \log(K/w_a)
 =\frac32\log(1/a)-\frac1q\log\ell_a
      +\frac12\log2+o(1)
 \sim\frac32\ell_a.
\]
Consequently, for every real $\gamma$,
\begin{equation}\label{eq:cube-ratio}
 \frac{X_a}{K^{2/3}w_a^{1/3}(1+\log(K/w_a))^\gamma}
 \sim 2^{-5/6}(3/2)^{-\gamma}
                   \ell_a^{\,2/(3q)-\gamma}.
\end{equation}
This tends to infinity for $\gamma<2/(3q)$ and to the positive constant
$2^{-5/6}(2/3)^{2/(3q)}$ at equality. The uniform-cube part of
Theorem~\ref{thm:main} follows.

\section{A fixed cube-truncated Gaussian}
Let $\rho$ be the standard Gaussian conditioned on $[-1,1]^d$.
Its density is the product of copies of
\[
 f(x)=Z_1^{-1}e^{-x^2/2}\ind_{[-1,1]}(x),\qquad
 Z_1=\int_{-1}^1 e^{-s^2/2}\dd s.
\]
Write its first two coordinates as independent symmetric variables
$X,Z$, and put $\ell=f(1)>0$. For $0\le s\le2$, define
\[
 F(s)=\Pr(X>1-s)=\int_{1-s}^1 f(x)\dd x.
\]
This function is continuous and strictly increasing. Smoothness of
$f$ up to the right endpoint from the interior gives
\begin{equation}\label{eq:F}
 F(s)=\ell s+O(s^2)\qquad(s\downarrow0),
\end{equation}
uniformly for $0\le s\le Ca$ at every fixed $C$.

Fix once and for all $k=1/8$. For small $a>0$, put $b=ka$ and
$m_a=F(a)$. Choose $c_a\in[(1-k)a,(1+k)a]$ by
\begin{equation}\label{eq:match}
 \E F(c_a+bZ)=F(a).
\end{equation}
Existence and uniqueness follow directly: throughout the candidate
interval, $c+bZ$ lies in $[(1-2k)a,(1+2k)a]\subset(0,2)$ for small
$a$. At its lower endpoint the widths are at most $a$, strictly
less almost surely, so the expectation is below $F(a)$; at the
upper endpoint it is above $F(a)$. Continuity and strict monotonicity
complete the argument. Since $\E Z=0$, \eqref{eq:F} yields
\begin{equation}\label{eq:ca}
 c_a=a+O(a^2),\qquad m_a\sim\ell a.
\end{equation}

Take
\[
 A_a=\{x_1>1-a\},\qquad
 B_a=\{x_1+b x_2>1-c_a\}.
\]
Their masses are exactly $m_a$ by \eqref{eq:match}. On each
$x_2=z$ slice their intervals are nested and have probabilities
$F(a)$ and $F(c_a+bz)$. Thus their symmetric-difference mass is exactly
\begin{equation}\label{eq:D}
 D_a=\E|F(c_a+bZ)-F(a)|.
\end{equation}
Equations \eqref{eq:F}--\eqref{eq:ca} imply, uniformly for
$z\in[-1,1]$,
\[
 F(c_a+bz)-F(a)=\ell k a z+O(a^2).
\]
The absolute-value function is $1$-Lipschitz, so taking absolute
values and expectations gives
\begin{equation}\label{eq:Dasymptotic}
 D_a=\ell k\E|Z|\,a+O(a^2),\qquad
 D_a/m_a\longrightarrow k\E|Z|>0.
\end{equation}
The overlap mass is exactly $m_a-D_a/2$.

Now choose
\[
 n_a=\frac{e_1+b e_2}{\sqrt{1+b^2}},\qquad
 \eta_a=|e_1-n_a|,\qquad
 R_a=K[\log(1/m_a)]^{1/q}.
\]
The convex potentials
\[
 u_a(x)=R_a(x_1-1+a)_+,\qquad
 v_a(x)=\frac{R_a}{\sqrt{1+b^2}}(x_1+b x_2-1+c_a)_+
\]
give the exact Brenier maps $R_a e_1\ind_{A_a}$ and
$R_a n_a\ind_{B_a}$, by the same direct supporting-affine-function
argument. Their target laws are
\[
 \mu_a=(1-m_a)\delta_0+m_a\delta_{R_ae_1},\qquad
 \nu_a=(1-m_a)\delta_0+m_a\delta_{R_an_a}.
\]
For small $a$, $m_a<1/2$ and $e_1\cdot n_a>0$. Lemma~\ref{lem:atoms}
proves that the pair has combined exponential moment $4-2m_a\le4$
and satisfies the exact formulas
\begin{equation}\label{eq:gaussian-exact}
 w_a^2=m_aR_a^2\eta_a^2,\qquad
 X_a^2=R_a^2[D_a+\eta_a^2(m_a-D_a/2)].
\end{equation}
Also
\[
 \eta_a^2=2-\frac2{\sqrt{1+k^2a^2}}\sim k^2a^2.
\]
Write $L_a=\log(1/m_a)\sim\log(1/a)$. Combining these formulas gives
\begin{align*}
 w_a&\sim Kk\sqrt\ell\,a^{3/2}L_a^{1/q}\longrightarrow0,\\
 X_a&\sim K\sqrt{\ell k\E|Z|}\,a^{1/2}L_a^{1/q},\\
 \log(K/w_a)&\sim\frac32 L_a.
\end{align*}
For the second relation, the angular overlap term in
\eqref{eq:gaussian-exact} is $O(a^3)$ inside the brackets, whereas
$D_a$ is asymptotic to a positive constant times $a$.
It follows that, for every real $\gamma$,
\begin{equation}\label{eq:gaussian-ratio}
 \frac{X_a}{K^{2/3}w_a^{1/3}(1+\log(K/w_a))^\gamma}
 \sim\ell^{1/3}k^{1/6}\sqrt{\E|Z|}\,(3/2)^{-\gamma}
                  L_a^{\,2/(3q)-\gamma}.
\end{equation}
This proves both the divergence for every smaller logarithmic power
and the finite positive limiting ratio at $\gamma=2/(3q)$. It
completes the lower-bound proof of Theorem~\ref{thm:main}.

\section{Upper bound and full Gaussian contrast}
Both fixed sources above satisfy manuscript 007, Corollary 3.2:
the uniform cube is compact log-concave, and the truncated Gaussian
is strongly log-concave with $\kappa=1$. Their slice variation is
finite. Its Corollary 4.1 therefore gives
\[
 X_\rho(\mu,\nu)
 \le C_{q,\rho,d}K^{2/3}w^{1/3}
           (1+\log(K/w))^{2/(3q)},\qquad 0<w\le K,
\]
on $\mathcal E_{q,K}$. This is the stated matching upper bound,
using the source-dependent potential theorems proved in manuscript
001 v3, Theorems 1.1 and 1.4. The present note does not independently
revalidate those theorems. In particular, subgaussian target tails ($q=2$)
have a necessary logarithmic power $1/3$ for either hard-boundary
source.

The conclusion differs for the full standard Gaussian $\gamma_d$.
Use target moment order $p=3$, independently of the stretched-exponential
parameter $q>0$. Since
\[
 s^3\le C_q e^{s^q},\qquad
 C_q=\sup_{s\ge0}s^3e^{-s^q}=(3/(qe))^{3/q},
\]
the class assumption bounds the sum of third moments by $4C_qK^3$.
Thus the sum $M$ of their third-moment roots satisfies
\[
 M\le2^{2/3}(4C_qK^3)^{1/3}.
\]
Manuscript 001 v3, Theorem 1.6, applied with moment order three,
gives the pure bound
\[
 \norm{T_\mu-T_\nu}_{L^2(\gamma_d)}
 \le C_{d,q}K^{2/3}W_2(\mu,\nu)^{1/3}
 \qquad ((\mu,\nu)\in\mathcal E_{q,K}),
\]
with no logarithmic factor. Its one-third exponent is sharp for
$d\ge2$: the three-atom bounded-target examples of Theorem 1.3 in
that manuscript can be scaled into the ball of radius
$K(\log2)^{1/q}$, which places their combined exponential moment
at most four. Thus a logarithmic loss is necessary for the fixed
hard-boundary strongly log-concave example but absent for the fixed
full Gaussian.

\begin{remark}[Dimension and other exponential bounds]
The lower constructions require $d\ge2$. In dimension one, an
atomless source has the exact monotone-quantile identity
$X_\rho=W_2$. The number four in the exponential-moment class is
only a convenient normalization. For any fixed $B>2$, replace
$R=K[\log(1/m)]^{1/q}$ by
$R=K[\log(c/m)]^{1/q}$ with $c=(B-2)/2$, for sufficiently small
$m<c$. The combined exponential moment becomes $B-2m\le B$,
and the same asymptotics and logarithmic exponent hold. At $B=2$
the class contains only the pair $(\delta_0,\delta_0)$.
\end{remark}

\begin{thebibliography}{9}
\raggedright
\bibitem{M001}
mxym/math research collection, \emph{Sharp Brenier stability under target moment bounds},
version 3, 7 October 2026, manuscript 001, \texttt{mxym/math}.
Theorem 1.7 and Section 11 give the cube-truncated Gaussian
two-atom family; Section 12.2 gives the exact uniform-cube family.
Theorem 1.6 gives the full-Gaussian finite-moment one-third bound;
Theorem 1.3 gives its bounded-target sharpness. Theorems 1.1 and 1.4
prove the source-dependent potential estimates used for the upper bounds.
Audited source blob
\texttt{0007edf5e28e4b4b0d34b2043cffd6e13e9f5aed}.
Pinned source:
\url{https://github.com/mxym/math/blob/0bd6c682edbe6f15001ddc2d73aac3c84ac979a4/preprints/001-strongly-log-concave-brenier/v3/manuscript.tex}.
\bibitem{M007}
mxym/math research collection, \emph{Tail-sensitive convex-gradient interpolation and
Brenier stability beyond bounded targets}, version 1,
7 October 2026, manuscript 007, \texttt{mxym/math}.
Theorem 3.1, Corollary 3.2, and Corollary 4.1.
Audited source blob \texttt{2d6f87539f93f87b33e2484334542ab924143057}.
Pinned source:
\url{https://github.com/mxym/math/blob/0bd6c682edbe6f15001ddc2d73aac3c84ac979a4/preprints/007-tail-brenier-stability/v1/main.tex}.
\end{thebibliography}
\noindent Both sources are read at repository commit
\texttt{0bd6c682edbe6f15001ddc2d73aac3c84ac979a4}.
The note records a new combination of their explicit results;
it makes no claim of novelty relative to the literature.
\end{document}
